\chapter{Hodge-theoretic preliminaries}
\label{chap:hodge}

% archon:covers MR4736527GeometricLocalSystemsOnVeryGeneralCurvesAndIsomonodromy/HodgeTheory.lean

\providecommand{\id}{\operatorname{id}}

We recall the definition of a complex variation of Hodge structure and of the
parabolic Higgs bundle associated to it, and we record the positivity property of
the top Hodge piece (\cref{lem:positive-Hodge-bundle}) together with its
instability consequence (\cref{cor:unstable}). These facts feed the parabolic
semistability arguments used elsewhere in the project.

\section{Complex variations of Hodge structure}

\begin{definition}[Polarizable complex variation of Hodge structure]
  \label{def:complex-variation}
  \lean{LandesmanLitt.ComplexVHS}
  \uses{def:flat-bundle}
  % SOURCE: [MR4736527 / arXiv:2202.00039], Definition (Polarizable complex
  % variations of Hodge structure), p. of §4.1
  % (read from references/MR4736527-local-systems-isomonodromy.tex, lines 1732-1751)
  % SOURCE QUOTE: "A {\em complex variation of Hodge structure} on $X$ is a triple $(V,
  %	V^{p,q}, D)$, where $V$ is a $C^\infty$ complex vector bundle on $X$,
  %	$V=\oplus V^{p,q}$ is a direct sum decomposition, and $D$ is a flat
  %	connection satisfying Griffiths transversality: $$D(V^{p,q})\subset
  %	A^{1,0}(V^{p,q})\oplus A^{0,1}(V^{p,q})\oplus A^{1,0}(V^{p-1,q+1})\oplus
  %	A^{0,1}(V^{p+1, q-1}).$$
  %	A {\em polarization} on $(V, V^{p,q}, D)$ is a flat Hermitian form $\psi$ on $V$ such that the $V^{p,q}$ are orthogonal to one another under $\psi$, and such that $(-1)^p\psi$ is positive definite on each $V^{p,q}$.
  %	A {\em polarizable complex variation of Hodge structure} is a complex
  %	variation of Hodge structure which admits a polarization."
  % SOURCE QUOTE (associated flat bundle and Hodge filtration): "We call the holomorphic flat vector bundle $(E,
  %    \nabla):=(\ker(D)\otimes_{\mathbb{C}} \mathscr{O}, \on{id}\otimes d)$ the
  %    {\em holomorphic flat vector bundle associated to the complex variation of
  %    Hodge structure}. The filtration $F^pV:=\oplus_{j\geq p} V^{j,q}$ of $V$ induces a decreasing \emph{Hodge filtration} $F^\bullet V$ by holomorphic sub-bundles, such that $\nabla(F^p)\subset F^{p-1}\otimes \Omega^1_X.$"
  \textit{Source: MR4736527 (arXiv:2202.00039), §4.1, Definition of polarizable
  complex variation of Hodge structure.}

  Let \(X\) be a smooth irreducible complex variety. A \emph{complex variation of
  Hodge structure} on \(X\) is a triple \((V, V^{p,q}, D)\) where \(V\) is a
  \(C^\infty\) complex vector bundle on \(X\), \(V = \bigoplus_{p,q} V^{p,q}\) is a
  direct sum decomposition into \(C^\infty\) complex sub-bundles, and \(D\) is a
  flat connection on \(V\) satisfying the \emph{Griffiths transversality}
  condition
  \[
    D(V^{p,q}) \subset A^{1,0}(V^{p,q}) \oplus A^{0,1}(V^{p,q})
    \oplus A^{1,0}(V^{p-1,q+1}) \oplus A^{0,1}(V^{p+1,q-1}),
  \]
  where \(A^{a,b}(W)\) denotes the smooth \((a,b)\)-forms valued in \(W\).

  A \emph{polarization} on \((V, V^{p,q}, D)\) is a flat Hermitian form \(\psi\) on
  \(V\) under which the \(V^{p,q}\) are mutually orthogonal and such that
  \((-1)^p \psi\) is positive definite on each \(V^{p,q}\). A \emph{polarizable}
  complex variation of Hodge structure is one which admits a polarization.

  The \emph{holomorphic flat vector bundle associated} to the variation is
  \((E, \nabla) := (\ker(D) \otimes_{\C} \sO,\ \id \otimes d)\), a flat bundle in
  the sense of \cref{def:flat-bundle}. The filtration \(F^p V := \bigoplus_{j \geq p}
  V^{j,q}\) induces a decreasing \emph{Hodge filtration} \(F^\bullet V\) by
  holomorphic sub-bundles satisfying
  \begin{equation}\label{eqn:griffiths-transversality}
    \nabla(F^p) \subset F^{p-1} \otimes \Omega^1_X.
  \end{equation}
  A complex local system \(\bV\) on \(X\) \emph{underlies a polarizable complex
  variation of Hodge structure} if \(\bV \simeq \ker(D)\) for some polarizable
  complex variation of Hodge structure \((V, V^{p,q}, D)\).
\end{definition}

\begin{definition}[Deligne canonical extension]
  \label{def:deligne-canonical-extension}
  \lean{LandesmanLitt.DeligneCanonicalExtension}
  \uses{def:flat-bundle}
  % SOURCE: [MR4736527 / arXiv:2202.00039], Definition (Deligne canonical
  % extension), §4.1
  % (read from references/MR4736527-local-systems-isomonodromy.tex, lines 1759-1769)
  % SOURCE QUOTE: "Let $\overline{X}$ be a smooth projective variety containing $X$ as a dense open subvariety with simple normal crossings complement $Z$.
  %    Let $(E, \nabla)$ be a flat holomorphic vector bundle on $X$. The \emph{Deligne canonical extension}
  %      $(\overline{E}, \overline{\nabla}: \overline{E}\to \overline{E}\otimes
  %      \Omega^1_{\overline{X}}(\log Z))$ of $(E, \nabla)$ to $\overline{X}$ is the
  %      unique flat vector bundle on $\overline{X}$ with regular singularities
  %      along $Z$, equipped with an isomorphism $(\overline{E},
  %      \overline{\nabla})|_X\overset{\sim}{\to}(E, \nabla)$, characterized by the
  %      property that all eigenvalues of its residues
  %      along components of $Z$
  %      have real parts lying in $[0,1)$."
  \textit{Source: MR4736527 (arXiv:2202.00039), §4.1, Definition of the Deligne
  canonical extension \cite[Remarques 5.5(i)]{deligne:regular-singular}.}

  Let \(\overline{X}\) be a smooth projective variety containing \(X\) as a dense
  open subvariety with simple normal crossings complement \(Z\), and let
  \((E, \nabla)\) be a flat holomorphic vector bundle on \(X\). The \emph{Deligne
  canonical extension}
  \[
    \bigl(\overline{E},\ \overline{\nabla} : \overline{E} \to
    \overline{E} \otimes \Omega^1_{\overline{X}}(\log Z)\bigr)
  \]
  of \((E, \nabla)\) to \(\overline{X}\) is the unique flat vector bundle on
  \(\overline{X}\) with regular singularities along \(Z\), equipped with an
  isomorphism \((\overline{E}, \overline{\nabla})|_X \xrightarrow{\sim} (E,
  \nabla)\), characterized by the property that all eigenvalues of its residues
  along the components of \(Z\) have real parts lying in \([0,1)\).
\end{definition}

\section{The associated parabolic Higgs bundle}

\begin{definition}[The associated parabolic Higgs bundle]
  \label{def:parabolic-higgs}
  \lean{LandesmanLitt.ParabolicHiggs}
  \uses{def:complex-variation, def:parabolic-bundle, def:deligne-canonical-extension}
  % SOURCE: [MR4736527 / arXiv:2202.00039], Definition (The associated Higgs
  % bundle), §4.1
  % (read from references/MR4736527-local-systems-isomonodromy.tex, lines 1772-1792)
  % SOURCE QUOTE: "Let $(E, F^\bullet, \nabla)$ be a holomorphic vector bundle $E$ on a smooth variety $\overline{X}$, with a flat connection $\nabla$ with regular singularities along a simple normal crossings divisor $Z\subset\overline{X}$, and a decreasing filtration $F^\bullet$ by holomorphic sub-bundles satisfying the Griffiths transversality condition $\nabla(F^p)\subset F^{p-1}\otimes \Omega^1_{\overline{X}}(\log Z).$ The \emph{associated Higgs bundle} is the pair $(\oplus_i \on{gr}^i_{F^\bullet}E, \theta)$, where the \emph{Higgs field} $$\theta:=\bigoplus_i (\theta_i: \on{gr}^i_{F^\bullet}E\to \on{gr}^{i-1}_{F^\bullet}E\otimes\Omega^1_{\overline{X}}(\log Z))$$ is the $\mathscr{O}_{\overline{X}}$-linear map induced by $\nabla.$"
  % SOURCE QUOTE (parabolic structure): "The vector bundle $E$ canonically has the structure of a parabolic bundle
  %    $E_\star$ ... This structure induces the structure of a parabolic bundle on $\oplus_i
  %    \on{gr}^i_{F^\bullet} E_\star$, ... as a direct sum of subquotients of $E$,
  %    preserved by $\theta$. We refer to the pair $(\oplus_i
  %    \on{gr}^i_{F^\bullet}E_\star, \theta)$ with its parabolic structure as the \emph{parabolic Higgs bundle} associated to the variation of Hodge structure."
  \textit{Source: MR4736527 (arXiv:2202.00039), §4.1, Definition of the
  associated Higgs bundle.}

  Let \((E, F^\bullet, \nabla)\) consist of a holomorphic vector bundle \(E\) on a
  smooth variety \(\overline{X}\), a flat connection \(\nabla\) with regular
  singularities along a simple normal crossings divisor \(Z \subset \overline{X}\),
  and a decreasing filtration \(F^\bullet\) by holomorphic sub-bundles satisfying
  the Griffiths transversality condition
  \begin{equation}\label{eqn:griffiths-transversality-2}
    \nabla(F^p) \subset F^{p-1} \otimes \Omega^1_{\overline{X}}(\log Z).
  \end{equation}
  The \emph{associated Higgs bundle} is the pair \((\bigoplus_i \gr^i_{F^\bullet} E,
  \theta)\), where the \emph{Higgs field}
  \[
    \theta := \bigoplus_i \bigl(\theta_i : \gr^i_{F^\bullet} E \to
    \gr^{i-1}_{F^\bullet} E \otimes \Omega^1_{\overline{X}}(\log Z)\bigr)
  \]
  is the \(\sO_{\overline{X}}\)-linear map induced by \(\nabla\) (the
  Kodaira--Spencer map); it is \(\sO\)-linear because Griffiths transversality
  forces \(\nabla\) to shift the filtration by exactly one step.

  The bundle \(E\) carries a canonical parabolic structure \(E_\star\) in the sense
  of \cref{def:parabolic-bundle}. This induces a parabolic structure on
  \(\bigoplus_i \gr^i_{F^\bullet} E_\star\), realized as a direct sum of
  subquotients of \(E\) and preserved by \(\theta\). The pair
  \((\bigoplus_i \gr^i_{F^\bullet} E_\star, \theta)\) with this parabolic structure
  is the \emph{parabolic Higgs bundle associated to the variation of Hodge
  structure}. When the input \((E, F^\bullet, \nabla)\) is the Deligne canonical
  extension of the flat bundle associated (\cref{def:complex-variation}) to a
  polarizable complex variation of Hodge structure, this is the parabolic Higgs
  bundle of that variation.
\end{definition}

\section{Basic facts}

We collect some basic facts about polarizable complex variations of Hodge
structure, the canonical extensions thereof, and their associated Higgs bundles.

\begin{proposition}[Basic facts on complex variations of Hodge structure]
  \label{prop:basic-facts}
  \lean{LandesmanLitt.basic_facts}
  \uses{def:complex-variation, def:parabolic-higgs, def:associated-parabolic, def:par-degree, def:parabolic-stability, def:deligne-canonical-extension}
  % SOURCE: [MR4736527 / arXiv:2202.00039], Proposition prop:basic-facts, §4.1
  % (read from references/MR4736527-local-systems-isomonodromy.tex, lines 1794-1830)
  % SOURCE QUOTE: "Let $\overline{X}$ be a smooth projective curve, $Z\subset
  % \overline{X}$ a reduced divisor, and let
  % $X=\overline{X}\setminus Z$.
  % Let $(V, V^{p,q}, D)$ be a polarizable complex variation of Hodge structure on
  % $X$,
  % and let $(E, F^\bullet,
  % \nabla)$ be the holomorphic flat vector bundle associated to this variation of
  % Hodge structure, with its Hodge filtration. Let $(\overline{E}, \overline{\nabla})$ be its Deligne canonical extension.
  % Let $\overline{E}_\star$ be the parabolic bundle associated to $(\overline{E}, \nabla)$, as defined in
  % \autoref{definition:associated-parabolic}.
  % \begin{enumerate}
  %     \item The local system $\mathbb{V}:=\ker(\nabla)$ associated to $(E, \nabla)$ is semisimple.
  %     \item The local system $\mathbb{V}$ may be canonically decomposed as $$\mathbb{V}\simeq \bigoplus_i \mathbb{L}_i\otimes
  % 	    W_i,$$ where the $\mathbb{L}_i$ are pairwise non-isomorphic
  % 	    irreducible complex local systems on $X$, and each $W_i$ is a complex vector
  % 	    space. Each $\mathbb{L}_i$ underlies a polarizable complex
  % 	    variation of Hodge structure, and each $W_i$ carries a complex polarized Hodge structure, both unique up to shifting the grading, and compatible with the variation carried by $\mathbb{V}$.
  %     \item $\overline{E}_\star$ has parabolic degree zero.
  %     \item There exists a canonical extension of $F^\bullet$ to $\overline{E}$, such that $(\overline{E}, F^\bullet, \overline{\nabla})$ satisfies the Griffiths transversality condition (\ref{eqn:griffiths-transversality-2}).
  %     \item The parabolic Higgs bundle $(\oplus_i
  % 	    \on{gr}^i_{F^\bullet}\overline{E}_\star, \theta)$ associated
  % 	    to $(\overline{E}_\star, F^\bullet, \overline{\nabla})$ is
  % 	    parabolically polystable of degree zero.
  % 	    That is, there exist a collection of parabolic vector bundles
  % 	    $E^j_\star$ with
  % 	    $\deg(E^j_\star) = 0$ and maps
  % 	    $\theta^j : E^j_\star \to E^j_\star \otimes \Omega^1_{\overline{X}}(\log Z)$ so that
  % 	    both
  % 	    $(\oplus_i\on{gr}^i_{F^\bullet}\overline{E}_\star,
  % 	    \theta)=\oplus_j (E^j_\star, \theta^j)$
  % 	    and for any
  % 	    $\theta^j$-stable proper sub-bundle $H_\star \subset
  % 	    E^j_\star$ with its induced parabolic structure, $\on{par-deg}
  % 	    H_\star < 0$."
  \textit{Source: MR4736527 (arXiv:2202.00039), §4.1, Proposition prop:basic-facts.}

  Let \(\overline{X}\) be a smooth projective curve, \(Z \subset \overline{X}\) a
  reduced divisor, and let \(X = \overline{X} \setminus Z\). Let
  \((V, V^{p,q}, D)\) be a polarizable complex variation of Hodge structure on
  \(X\), and let \((E, F^\bullet, \nabla)\) be the holomorphic flat vector bundle
  associated to this variation, with its Hodge filtration. Let
  \((\overline{E}, \overline{\nabla})\) be its Deligne canonical extension, and let
  \(\overline{E}_\star\) be the parabolic bundle associated to
  \((\overline{E}, \nabla)\) in the sense of \cref{def:associated-parabolic}. Then:
  \begin{enumerate}
    \item the local system \(\bV := \ker(\nabla)\) associated to \((E, \nabla)\) is
      semisimple;
    \item \(\bV\) decomposes canonically as
      \(\bV \simeq \bigoplus_i \bL_i \otimes W_i\), where the \(\bL_i\) are pairwise
      non-isomorphic irreducible complex local systems on \(X\) and each \(W_i\) is
      a complex vector space; each \(\bL_i\) underlies a polarizable complex
      variation of Hodge structure and each \(W_i\) carries a complex polarized
      Hodge structure, both unique up to shifting the grading and compatible with
      the variation carried by \(\bV\);
    \item \(\overline{E}_\star\) has parabolic degree zero;
    \item there exists a canonical extension of \(F^\bullet\) to \(\overline{E}\)
      such that \((\overline{E}, F^\bullet, \overline{\nabla})\) satisfies the
      Griffiths transversality condition \eqref{eqn:griffiths-transversality-2};
    \item the associated parabolic Higgs bundle
      \((\bigoplus_i \gr^i_{F^\bullet} \overline{E}_\star, \theta)\) of
      \((\overline{E}_\star, F^\bullet, \overline{\nabla})\) is parabolically
      polystable of degree zero. That is, there exist parabolic vector bundles
      \(E^j_\star\) with \(\deg(E^j_\star) = 0\) and maps
      \(\theta^j : E^j_\star \to E^j_\star \otimes \Omega^1_{\overline{X}}(\log Z)\)
      such that
      \((\bigoplus_i \gr^i_{F^\bullet} \overline{E}_\star, \theta) =
      \bigoplus_j (E^j_\star, \theta^j)\), and for every \(\theta^j\)-stable proper
      sub-bundle \(H_\star \subset E^j_\star\) with its induced parabolic structure
      one has \(\operatorname{par-deg} H_\star < 0\).
  \end{enumerate}
\end{proposition}

% SOURCE QUOTE PROOF: "The proof of
% (1) is explained in \cite[1.11-1.12]{deligne1987theoreme} and (2) is
% \cite[1.13]{deligne1987theoreme}. The proof of (3) follows from
% \cite[B.3]{esnault1986logarithmic}
% while (4) is explained in e.g.~\cite[Section 7]{brunebarbe2017semi}.
% Finally,
% (5) is due to Simpson \cite[Theorem 5]{simpson1990harmonic}.
% See the discussion in the introduction of \cite{arapura2019vanishing} for a nice summary of this and related topics."
\begin{proof}
  This is an external, black-box result; we record the literature references the
  paper gives. Parts (1) and (2) are explained in \cite[1.11--1.12]{deligne1987theoreme}
  and \cite[1.13]{deligne1987theoreme} respectively. Part (3) follows from
  \cite[B.3]{esnault1986logarithmic}, and part (4) is explained in
  e.g.\ \cite[Section 7]{brunebarbe2017semi}. Finally, part (5) is due to Simpson
  \cite[Theorem 5]{simpson1990harmonic}. See the introduction of
  \cite{arapura2019vanishing} for a summary of this and related topics.
\end{proof}

\section{Positivity and instability}

The next lemma shows that the top filtered piece of the Hodge filtration tends to
have positive parabolic degree, the key input to the semistability arguments used
elsewhere in the project.

\begin{lemma}[Positivity of the top Hodge bundle]
  \label{lem:positive-Hodge-bundle}
  \lean{LandesmanLitt.positive_Hodge_bundle}
  \uses{def:complex-variation, def:parabolic-higgs, def:par-slope, def:deligne-canonical-extension}
  % SOURCE: [MR4736527 / arXiv:2202.00039], Lemma 4.? (positive-Hodge-bundle),
  % §4.1
  % (read from references/MR4736527-local-systems-isomonodromy.tex, lines 1848-1859)
  % SOURCE QUOTE: "Let $C$ be a smooth proper curve, $Z\subset C$ a reduced effective divisor,
  %    and $(V, V^{p,q}, D)$ a  a polarizable complex variation of Hodge structure
  %    on $C\setminus Z$. Let $(\overline{E}_\star, F^\bullet, \nabla)$ be the Deligne canonical
  %    extension of the associated flat holomorphic vector bundle to $C$ with its
  %    canonical parabolic structure.  Let $i$ be maximal such that
  %    $F^i\overline{E}_\star$ is non-trivial, where $F^\bullet$ is the Hodge
  %    filtration, and suppose that the Higgs field $$\theta_i:
  %    F^i\overline{E}_\star\to
  %    \on{gr}^{i-1}_{F^{\bullet}}\overline{E}_\star\otimes \Omega^1_C(\log Z)$$ is
  %    non-zero. Then $F^i\overline{E}_\star$ has positive parabolic degree."
  \textit{Source: MR4736527 (arXiv:2202.00039), §4.1, Lemma on positivity of the
  Hodge bundle.}

  Let \(C\) be a smooth proper curve, \(Z \subset C\) a reduced effective divisor,
  and \((V, V^{p,q}, D)\) a polarizable complex variation of Hodge structure on
  \(C \setminus Z\). Let \((\overline{E}_\star, F^\bullet, \nabla)\) be the Deligne
  canonical extension of the associated flat holomorphic vector bundle to \(C\),
  with its canonical parabolic structure. Let \(i\) be maximal such that
  \(F^i \overline{E}_\star\) is non-trivial, where \(F^\bullet\) is the Hodge
  filtration, and suppose the Higgs field
  \[
    \theta_i : F^i \overline{E}_\star \to
    \gr^{i-1}_{F^\bullet} \overline{E}_\star \otimes \Omega^1_C(\log Z)
  \]
  is non-zero. Then \(F^i \overline{E}_\star\) has positive parabolic degree.
\end{lemma}

% SOURCE QUOTE PROOF: "The parabolic vector bundle $F^i\overline{E}_\star$ with the zero Higgs field is a quotient of
% the parabolic Higgs bundle  $(\oplus_i \on{gr}^i_{F^\bullet}\overline{E}_\star, \theta)$, and
% hence has non-negative parabolic degree by the fact that the latter is polystable, hence semistable, of
% degree zero, by \autoref{prop:basic-facts}(5). It has degree zero if and only if
% it is a direct summand of $(\oplus_i \on{gr}^i_{F^\bullet}\overline{E}_\star, \theta)$ by polystability; but this is ruled out by the nonvanishing of $\theta_i$."
\begin{proof}
  \uses{prop:basic-facts, def:parabolic-stability, def:par-slope}
  The parabolic vector bundle \(F^i \overline{E}_\star\), equipped with the zero
  Higgs field, is a quotient of the associated parabolic Higgs bundle
  \((\bigoplus_i \gr^i_{F^\bullet} \overline{E}_\star, \theta)\) of
  \cref{def:parabolic-higgs}. By \cref{prop:basic-facts}(5) that parabolic Higgs
  bundle is parabolically polystable, hence semistable, of parabolic degree zero.
  Therefore the quotient \(F^i \overline{E}_\star\) has non-negative parabolic
  degree.

  Its parabolic degree is zero if and only if \(F^i \overline{E}_\star\) is a
  direct summand of the polystable parabolic Higgs bundle as a Higgs sub-object.
  The non-vanishing of \(\theta_i\) means that \(F^i \overline{E}_\star\) with its
  zero Higgs field is not preserved as a Higgs sub-object, so it cannot be such a
  direct summand. Hence its parabolic degree is strictly positive.
\end{proof}

\begin{remark}
  \label{rem:positive-Hodge-bundle-attribution}
  This lemma is essentially due to Griffiths in the case of real variations of
  Hodge structure on a smooth proper curve, where it follows from the curvature
  formula for Hodge bundles. The statement above, for complex variations on a
  quasi-projective curve, is recorded with the short proof above because the
  authors found no precise reference in that generality.
\end{remark}

We now spell out how \cref{lem:positive-Hodge-bundle} relates to semistability.

\begin{corollary}[Instability of the parabolic bundle]
  \label{cor:unstable}
  \lean{LandesmanLitt.unstable}
  \uses{lem:positive-Hodge-bundle, def:parabolic-higgs, def:complex-variation}
  % SOURCE: [MR4736527 / arXiv:2202.00039], Corollary 4.? (unstable), §4.1
  % (read from references/MR4736527-local-systems-isomonodromy.tex, lines 1893-1897)
  % SOURCE QUOTE: "Let $(\overline{E}_\star, F^\bullet, \nabla)$ be as in
  % \autoref{lemma:positive-Hodge-bundle}. Then the parabolic bundle
  % $\overline{E}_\star$ is not semistable."
  \textit{Source: MR4736527 (arXiv:2202.00039), §4.1, Corollary on instability.}

  Let \((\overline{E}_\star, F^\bullet, \nabla)\) be as in
  \cref{lem:positive-Hodge-bundle}. Then the parabolic bundle
  \(\overline{E}_\star\) is not semistable. Equivalently, a flat bundle underlying
  a polarizable complex variation of Hodge structure with more than one non-zero
  Hodge piece is parabolically unstable: parabolic semistability forces the Hodge
  filtration to have a single part.
\end{corollary}

% SOURCE QUOTE PROOF: "The parabolic vector bundle $\overline{E}_\star$ has parabolic degree zero
%    by \autoref{prop:basic-facts}(3). But by
%    \autoref{lemma:positive-Hodge-bundle}, $F^i\overline{E}_\star$ has positive degree, and is hence a destabilizing subsheaf."
\begin{proof}
  \uses{prop:basic-facts, def:parabolic-stability, def:par-slope}
  The parabolic vector bundle \(\overline{E}_\star\) has parabolic degree zero by
  \cref{prop:basic-facts}(3) (a basic property of the Deligne canonical extension
  of the flat bundle of a polarizable complex variation of Hodge structure). By
  \cref{lem:positive-Hodge-bundle}, the sub-bundle \(F^i \overline{E}_\star\) has
  positive parabolic degree. A parabolic sub-bundle of positive parabolic degree
  inside a parabolic bundle of parabolic degree zero is destabilizing, so
  \(\overline{E}_\star\) is not semistable.
\end{proof}
